\section{VTAB: cómo convertir lo «general» en un experimento reproducible}

Sin un protocol, «general» se vuelve fácilmente un término de mercadotecnia. \term{VTAB} (Visual Task Adaptation Benchmark) divide la evaluación en upstream pretraining, un adaptation budget fijo y varias clases de downstream tasks. Así, los investigadores pueden comparar la capacidad de transferencia de una representación en una amplia variedad de tareas y, al mismo tiempo, se ponen al descubierto los atajos que solo funcionan en un dominio de datos concreto.

\subsection{Upstream: primero se fija el problema, luego se comparan las representaciones}

En el primer paso de VTAB, los investigadores pueden aprender representaciones con sus propios datos y modelos. El benchmark no impone un upstream dataset ni una architecture, pero exige que la adaptación posterior se haga con las mismas tareas y el mismo presupuesto. De este modo, lo que se compara es «la transferibilidad que produce todo el esquema de preentrenamiento», no solo el número de parámetros de un backbone.

\lecturefigure{slide-07-vtab-upstream.jpg}{La etapa upstream de VTAB recibe los datos y el modelo de preentrenamiento, y entrega la representación que se va a evaluar.}{Video oficial 04:40; Zhai et al. 2019.}

\begin{importantbox}{Definición de upstream / downstream}
\term{upstream} se refiere a la etapa de preentrenamiento: se eligen los datos $D_u$, el objetivo $\mathcal{L}_u$ y los parámetros $\theta$. \term{downstream} se refiere a la etapa de la tarea objetivo: dado un conjunto etiquetado pequeño $D_t$, se aprende la cabeza de la tarea o se hace fine-tuning de los parámetros. En forma simplificada:
\[
\theta^{\star}=\arg\min_{\theta}\mathcal{L}_u(D_u;\theta),\qquad
\phi_t^{\star}=\arg\min_{\phi_t}\mathcal{L}_t(D_t;\theta^{\star},\phi_t).
\]
Que la segunda expresión permita o no actualizar $\theta^{\star}$ depende del protocol: linear probe, partial fine-tuning o full fine-tuning.
\end{importantbox}

\subsection{Task families: ningún domain debe dominar las conclusiones}

Un solo promedio oculta las diferencias entre tipos de tareas. VTAB divide las tareas en grupos como natural, specialized y structured. Las imágenes naturales se parecen a la distribución habitual del preentrenamiento. Las imágenes especializadas incluyen percepción remota, imágenes médicas o sensores específicos. Las tareas estructuradas se centran en la posición, el conteo, la profundidad y las relaciones. Una representación general debe ser estable en varios grupos, no destacar solo en los grupos semánticamente cercanos.

\lecturefigure{slide-08-vtab-task-families.jpg}{VTAB muestrea natural, specialized y structured tasks de un espacio amplio de tareas visuales.}{Video oficial 06:00.}

\begin{knowledgebox}{Lectura de la figura: antes del promedio, revisar las curvas por grupo}
Las natural tasks pueden premiar la transferencia de semántica y texturas. Las specialized tasks ponen a prueba el domain shift. Las structured tasks dependen más de las relaciones espaciales. Si un método solo aventaja ampliamente en el grupo natural, conviene ser prudente antes de llamarlo «general». Se recomienda reportar el macro average, el promedio de cada grupo, la peor tarea, la varianza y el adaptation cost, en lugar de dar una sola puntuación total.
\end{knowledgebox}

\subsection{Adaptation: comprobar si la representación sirve con datos limitados}

El último paso del benchmark consiste en adaptar y evaluar con un presupuesto fijo de ejemplos downstream. Lo que se compara aquí no es «quién alcanza la mayor exactitud tras un fine-tuning ilimitado», sino si la representación ya codificó estructura útil en un sistema de coordenadas que los datos escasos puedan aprovechar con facilidad. Al leer un transfer protocol, hay que registrar a la vez el número de ejemplos, los parámetros entrenables, el presupuesto de ajuste de hiperparámetros y la forma de la evaluación final, porque relajar cualquiera de ellos puede convertir «una mejor representación» en «un entrenamiento downstream más fuerte».

\lecturefigure{slide-09-vtab-transfer-protocol.jpg}{VTAB transfiere la upstream representation a un conjunto de tareas downstream y luego agrega el desempeño de la adaptación.}{Video oficial 06:44.}

\begin{knowledgebox}{Lectura de la figura: el adaptador también es una variable experimental}
Un linear probe mide la separabilidad lineal de la representación. El full fine-tuning mide si la inicialización es favorable. Un parameter-efficient adapter mide la plasticidad a bajo costo. Cada protocolo responde una pregunta distinta. Si A usa full fine-tuning y B usa linear probe, la diferencia de puntuación puede deberse sobre todo a la capacidad de adaptación y no a la upstream representation.
\end{knowledgebox}

\begin{warningbox}{El diseño del benchmark también tiene blind spots}
VTAB aporta amplitud, pero sigue siendo un conjunto finito de tareas. No cubre por completo detection, segmentation, video, open-vocabulary reasoning, calibration ni la seguridad. Además, tras años de optimizar contra un benchmark aparece el meta-overfitting: la comunidad de investigación convierte de forma indirecta la test suite en un objetivo de entrenamiento. Por eso, la conclusión debe redactarse como «más transferible bajo esta suite y este protocol».
\end{warningbox}

\teachervoice{El ponente describe VTAB como la versión formalizada del pequeño juego de pocos ejemplos presentado antes. Lo que el docente enfatiza no es que «19 tareas equivalen a toda la visión», sino que un protocol común obliga primero a los investigadores a responder en qué tipos de tareas es transferible la representación y cuál es el costo de adaptación.}

\subsection{Resumen de la sección}

VTAB descompone la representación general en upstream, task families y adaptation protocol. Se acerca más que la ImageNet accuracy por sí sola a la idea de «aprender rápidamente nuevas tareas visuales», pero sigue siendo necesario reportar por grupos, controlar la capacidad de adaptación y estar atentos al meta-overfitting del benchmark.
